\subsection{COMPLETION OF SUBSPACES AND PRODUCT SPACES}

PROPOSITION 18. Let X be a set, let \((\mathbf{Y}_{\lambda})_{\lambda \in \mathbf{L}}\) be a family of uniform spaces, and for each \(\lambda \in \mathbf{L}\) let \(f_{\lambda}\) be a mapping of X into \(\mathbf{Y}_{\lambda}\) . Let X carry the coarsest uniformity \(\mathcal{U}\) which makes all the \(f_{\lambda}\) uniformly continuous. Then the uniformity of the Hausdorff completion \(\hat{\mathbf{X}}\) of X is the coarsest for which all the mappings \(\hat{f}_{\lambda} \colon \hat{\mathbf{X}} \to \hat{\mathbf{Y}}_{\lambda} (\lambda \in \mathbf{L})\) (no. 7, Proposition 15) are uniformly continuous. Furthermore, if \(j_{\lambda}\) is the canonical mapping of \(\mathbf{Y}_{\lambda}\) into \(\hat{\mathbf{Y}}_{\lambda}\) , and if \(g_{\lambda} = j_{\lambda} \circ f_{\lambda}\) , then \(\hat{\mathbf{X}}\) may be identified with the closure in \(\prod_{\lambda \in \mathbf{L}} \hat{\mathbf{Y}}_{\lambda}\) of the image of X under the mapping \(x \to (g_{\lambda}(x))\) .

Let \(X'\) (resp. \(Y_{\lambda}'\) ) be the Hausdorff uniform space associated with \(X\) (resp. \(Y_{\lambda}\) ), and let \(f_{\lambda}' \colon X' \to Y_{\lambda}'\) be the uniformly continuous mapping which makes the diagram

\[
\begin{array}{c} \mathbf {X} \xrightarrow {f _ {\lambda}} \mathbf {Y} _ {\lambda} \\ i \downarrow \\ \mathbf {X ^ {\prime}} \xrightarrow [ f _ {j} ^ {I} ]{} \mathbf {Y} _ {\lambda} ^ {\prime} \end{array}
\]

commutative (i being the canonical mapping).

The transitivity of initial uniformities (§ 2, no. 3, Proposition 5) shows on the one hand that \(\mathcal{U}\) is the coarsest uniformity for which the mappings \(j_{\lambda} \circ f_{\lambda}: \mathbf{X} \to \mathbf{Y}_{\lambda}'\) are uniformly continuous, and on the other hand that \(\mathcal{U}\) is also the inverse image under \(i\) of the coarsest uniformity \(\mathcal{U}'\) on the set \(\mathbf{X}'\) for which the \(f_{\lambda}'\) are uniformly continuous. Now \(\mathcal{U}'\) is Hausdorff, for if \(x_{1}, x_{2}\) are two points of \(\mathbf{X}\) such that \(j_{\lambda}(f_{\lambda}(x_{1})) = j_{\lambda}(f_{\lambda}(x_{2}))\) for each \(\lambda \in L\) , then \((x_{1}, x_{2})\) belongs to all the entourages of \(\mathcal{U}\) and hence \(i(x_{1}) = i(x_{2})\) . Proposition 17 of no. 8 therefore shows that \(\mathcal{U}'\) is the uniformity of the Hausdorff space \(\mathbf{X}'\) associated with \(\mathbf{X}\) .

This being so, the bijection \(x' \to (f_{\lambda}'(x'))\) identifies X with a uniform subspace of the product \(\prod_{\lambda} Y_{\lambda}'\) ( \(\S\) 2, no. 6, Proposition 8). Since the \(Y_{\lambda}'\) are Hausdorff, each \(Y_{\lambda}'\) can be identified with a dense subspace of its completion \(\hat{Y}_{\lambda}\) , and hence \(\prod_{\lambda} Y_{\lambda}'\) can be identified with a dense subspace of \(\prod_{\lambda} \hat{Y}_{\lambda}\) (Chapter I, § 4, no. 3, Proposition 7). But \(\prod_{\lambda} \hat{Y}_{\lambda}\) is Hausdorff and complete (no. 5, Proposition 10); the closure \(\overline{X}'\) of \(X'\) in \(\prod_{\lambda} \hat{Y}_{\lambda}\) is therefore a complete Hausdorff subspace (no. 4, Proposition 8) which can be identified with the Hausdorff completion \(\hat{X}\) of X; under this identification the mappings \(\hat{f}_{\lambda}\) become the projections onto the factors \(\hat{Y}_{\lambda}\) , and the proposition is proved.

COROLLARY I. Let X be a uniform space and let i be the canonical mapping of X into its Hausdorff completion \(\hat{X}\); let A be a subspace of X and j: A → X the canonical injection. Then j: \(\hat{A} \rightarrow \hat{X}\) is an isomorphism of \(\hat{A}\) onto the closure of i(A) in \(\hat{X}\) .

COROLLARY 2. Let \((\mathbf{Y}_{\lambda})_{\lambda \in \mathbf{L}}\) be a family of uniform spaces. Then the Hausdorff completion of the product space \(\prod_{\lambda \in \mathbf{L}}\mathbf{Y}_{\lambda}\) is canonically isomorphic to the product \(\prod_{\lambda \in \mathbf{L}}\hat{\mathbf{Y}}_{\lambda}\) .

